\chapter*{本卷结构}
\addcontentsline{toc}{chapter}{本卷结构}

本卷按第三册原书的四个大块组织。每一章都对应原书的一组主题文件，并把概念落到同一条实践主线：\code{linux\_cpu\_inference}。

第一部分从模型资产到量化 kernel。你会运行 tiny MLP，检查 tokenizer 和 safetensors manifest，手算张量地址，比较 scalar、packed 和 AVX2 路径，并解释 int8 scale、accumulator 和 requantize 的边界。

第二部分进入 Transformer、KV cache、sampler 和服务运行时。你会跑 reference decoder trace，解释每个 checkpoint 的 shape、stride、dtype 和 layout，构造 KV cache 状态表，分析 admission、queue wait、TTFT、TPOT、取消和拒绝。

第三部分处理模型导入、真实模型加载闭环、memory planner 和 IR lowering。你会把 manifest、shape verifier、tensor role、state edge、workspace、liveness 和 executable plan 写成可检查报告。

第四部分收束到后端优化和验收。你会比较 CPU kernel 证据、GPU 后端边界、工业专项能力和回归门禁。最终交付不是“跑通一次”，而是可复查的工程档案：命令、输入、输出、trace、benchmark、不可验证边界和回归规则。
